\exercice*

\begin{enumerate}
  (* for i in ordre *)
    (* if i == 0 *)
      \item On calcule le taux d'évolution avec la formule :
        \[
          \frac{\text{Valeur finale} - \text{Valeur initiale}}{\text{Valeur initiale}}=\frac{(( finale0 )) - (( initiale0 ))}{(( initiale0 ))}=\frac{(( finale0 - initiale0 ))}{(( initiale0 ))}=(( (taux0 / 100)|facteur("2") ))=(( taux0 ))\%
        \]
        (* if contexte == 0 *)
          Le chiffre d'affaires a donc (* if taux0 > 0 *)augmenté(* else *)diminué(* endif *) de (( taux0|abs ))\%.
        (* elif contexte == 1 *)
          La production hebdomadaire de vélos a donc (* if taux0 > 0 *)augmenté(* else *)diminué(* endif *) de (( taux0|abs ))\%.
        (* elif contexte == 2 *)
          La population de loups a donc (* if taux0 > 0 *)augmenté(* else *)diminué(* endif *) de (( taux0|abs ))\%.
        (* elif contexte == 3 *)
          Le nombre d'ouvrages a donc (* if taux0 > 0 *)augmenté(* else *)diminué(* endif *) de (( taux0|abs ))\%.
        (* else *)
          Le nombre d'adhérents a donc (* if taux0 > 0 *)augmenté(* else *)diminué(* endif *) de (( taux0|abs ))\%.
        (* endif *)
    (* elif i == 1 *)
      \item
          (* if taux1 > 0 *)Augmenter(* else *)Diminuer(* endif *) une valeur de (( taux1|abs ))\% revient à la multiplier par $1(( taux1 | signe ))\frac{(( taux1|abs ))}{100}=1(( (taux1/100)|facteur("s") ))=(( mult1|facteur ))$.
          Donc la nouvelle valeur est $(( initiale1 ))\times (( mult1|facteur ))=(( finale1 ))$.
          (* if contexte == 0 *)
          L'année suivante, le chiffre d'affaires était donc de (( finale1 )) milliers d'euros.
          (* elif contexte == 1 *)
          L'année suivante, l'usine produisait donc (( finale1 )) vélos par semaine.
          (* elif contexte == 2 *)
          Il y avait l'année suivante (( finale1 )) loups dans le parc.
          (* elif contexte == 3 *)
          Il y avait l'année suivante (( finale1 )) milliers d'ouvrages dans la bibliothèque.
          (* elif contexte == 4 *)
          Il y avait l'année suivante (( finale1 )) adhérents au club de basket.
          (* endif *)
    (* else *)
      \item Complétons le schéma suivant.
        \begin{center}
          \begin{tikzpicture}
            \node[draw=black, circle, inner sep=0pt, minimum size=2em] (initial) at (0, 0) {?};
            \node[draw=black, circle, inner sep=0pt, minimum size=2em] (final) at (3, 0) {$(( finale2 ))$};
            \draw[-Stealth, bend left] (initial) to node[above] {$(( taux2|facteur("s") ))\%$} node[below]{$\times?$} (final);
            \draw[-Stealth, bend left] (final)   to node[below]{$?$} (initial);
          \end{tikzpicture}
        \end{center}
        La valeur initiale est inconnue, mais après une (* if taux2 > 0 *)augmentation(* else *)diminution(* endif *) de (( taux2|abs ))\%, la valeur finale est (( finale2 )). Or, (* if taux2 > 0 *)augmenter(* else *)diminuer(* endif *) une valeur de (( taux2|abs ))\% revient à la multiplier par $1(( taux2|signe ))\frac{(( taux2|abs ))}{100}=1(( (taux2/100)|facteur("s") ))=(( mult2|facteur ))$. En d'autres termes, pour passer de la valeur initiale à la valeur finale, on l'a multipliée par (( mult2|facteur )). Pour revenir à la valeur initiale, il faut donc diviser la valeur finale par (( mult2|facteur )), ce qui donne :
        (* if taux2 == -25 *)
        \[
        \frac{(( finale2 ))}{(( mult2|facteur ))}
        =\frac{(( finale2 ))}{3/4}
        =(( finale2 ))\times\frac{4}{3}
        =\frac{(( finale2 ))}{3}\times4
        =(( (finale2/3)|facteur ))\times4
        =(( initiale2 ))
        \]
        (* else *)
        \[
        \frac{(( finale2 ))}{(( mult2|facteur ))}
        =\frac{(( finale2 * 10 ))}{(( (mult2*10)|facteur ))}
        =\frac{(( finale2 // 10 )) \times 100}{(( (mult2*10)|facteur ))}
        =\frac{\cancel{(( (mult2 * 10)|facteur ))}\times (( (finale2 // (100 * mult2))|facteur )) \times 100}{\cancel{(( (mult2*10)|facteur ))}}
        =(( initiale2 ))
        \]
        (* endif *)
    
        \begin{center}
          \begin{tikzpicture}
            \node[draw=black, circle, inner sep=0pt, minimum size=2em] (initial) at (0, 0) {$(( initiale2 ))$};
            \node[draw=black, circle, inner sep=0pt, minimum size=2em] (final) at (3, 0) {$(( finale2 ))$};
            \draw[-Stealth, bend left] (initial) to node[above] {$(( taux2|facteur("s") ))\%$} node[below]{$\times(( mult2|facteur ))$} (final);
            \draw[-Stealth, bend left] (final)   to node[below]{$\div(( mult2|facteur ))$} (initial);
          \end{tikzpicture}
        \end{center}
        (* if contexte == 0 *)
          En (( annee2 )), le chiffre d'affaires de l'entreprise était donc de (( initiale2 )) milliers d'euros.
        (* elif contexte == 1 *)
          En (( annee2 )), l'usine a donc produit (( initiale2 )) vélos par semaine.
        (* elif contexte == 2 *)
          En (( annee2 )), il y avait donc (( initiale2 )) loups dans le parc.
        (* elif contexte == 3 *)
          En (( annee2 )), la bibliothèque possédait donc (( initiale2 )) milliers d'ouvrages.
        (* else *)
          En (( annee2 )), le club réunissait donc (( initiale2 )) adhérents.
        (* endif *)
    (* endif *)
  (* endfor *)
\end{enumerate}
